% Flujo del programa del ESP32. Arriba, el ciclo de lectura (núcleo 1); las cajas
% punteadas corren en paralelo (la tarea de red en el núcleo 0).
\begin{tikzpicture}[
  font=\footnotesize,
  inicio/.style={rectangle, rounded corners=3.5mm, draw=black!70, fill=black!8, thick,
    minimum width=2cm, minimum height=0.7cm},
  proceso/.style={caja=cSensor, minimum width=2.8cm, minimum height=0.8cm},
  accion/.style={caja=cAct, minimum width=2.8cm, minimum height=0.8cm},
  aviso/.style={caja=cProc, minimum width=2.4cm, minimum height=0.8cm},
  red/.style={caja=cApp, minimum width=2.8cm, minimum height=0.8cm},
  decision/.style={diamond, aspect=2, draw=cProc, fill=cProc!10, thick,
    align=center, inner sep=1pt},
  si/.style={font=\scriptsize, text=black!70},
]
  \node[inicio]   (start) at (0,0)     {Inicio};
  \node[proceso]  (conf)  at (0,-1.2)  {Bloquear la puerta,\\crear colas y red};
  \node[proceso]  (leer)  at (0,-2.6)  {Leer sensores};
  \node[decision] (d1)    at (0,-4.1)  {¿Sensor activo\\y armado?};
  \node[decision] (d2)    at (0,-5.9)  {¿Otro sensor\\en 10\,s?};
  \node[accion]   (crit)  at (0,-7.6)  {Buzzer 30\,s y\\bloquear puerta};
  \node[red]      (pub)   at (0,-9.1)  {Poner el evento\\en la cola};
  \node[aviso]    (leve)  at (3.6,-5.9) {Alerta leve};

  % Procesos paralelos
  \node[red, dashed]    (lat) at (3.8,0.3)  {Red (núcleo 0):\\publica la cola y\\un latido cada 30\,s};
  \node[accion, dashed] (cmd) at (3.8,-1.35) {Orden de abrir:\\desbloquear};
  \node[accion, dashed] (gas) at (3.8,-2.75) {Humo en el MQ-2:\\buzzer y desbloquear};
  \node[accion, dashed] (ab)  at (3.8,-4.15) {Puerta abierta 1\,min:\\buzzer hasta cerrar};

  \draw[flecha] (start) -- (conf);
  \draw[flecha] (conf) -- (leer);
  \draw[flecha] (leer) -- (d1);
  \draw[flecha] (d1) -- node[si, right] {Sí} (d2);
  \draw[flecha] (d2) -- node[si, right] {Sí} (crit);
  \draw[flecha] (crit) -- (pub);
  % Sin activación: vuelve a leer
  \draw[flecha] (d1.west) -- node[si, above] {No} ++(-0.55,0) |- (leer.west);
  % Un solo sensor: alerta leve y se publica
  \draw[flecha] (d2.east) -- node[si, above] {No} (leve.west);
  \draw[flecha] (leve.south) |- (pub.east);
  % Después de poner el evento en la cola, vuelve a leer
  \draw[flecha, draw=black!45] (pub.south) -- ++(0,-0.4) -| (-2.9,-2.6) -- (leer.west);
\end{tikzpicture}
